% =========================================================================
% RELATÓRIO DE INICIAÇÃO CIENTÍFICA (COM CAPA INSTITUCIONAL UFABC)
% =========================================================================
\documentclass[12pt,a4paper]{article}

\usepackage{geometry}
\geometry{
    a4paper,
    margin=2cm,
    headheight=15pt
}

\RequirePackage{fix-cm}

% =========================================================================
% CABEÇALHO UNIFICADO E DESIGN SYSTEM (NETWORK FLOW)
% =========================================================================

% =========================================================================
% PACOTES FUNDAMENTAIS
% =========================================================================
\usepackage{amssymb}
\usepackage{amsmath}
\usepackage{amsthm}
\usepackage{bm}
\usepackage{enumitem}
\usepackage{mathtools}
\usepackage{lipsum}
\usepackage[brazilian]{babel}
\usepackage[utf8]{inputenc}
\usepackage[T1]{fontenc}
\usepackage[pdftex]{graphicx}
\usepackage{color}
\usepackage{xcolor}
\usepackage{colortbl}
\usepackage{tikz}
\usetikzlibrary{arrows.meta, calc, positioning, decorations.markings}
\usepackage{verbatim}
\usepackage[hyphens]{url}
\usepackage{longtable}
\usepackage{booktabs}
\usepackage[section]{placeins}
\usepackage{listings}
\usepackage{anyfontsize}
\usepackage{lmodern}
\usepackage{helvet}
\renewcommand{\familydefault}{\sfdefault}
\usepackage{fancyhdr}

\newif\ifapendiceteorico
\apendiceteoricofalse
\newif\ifapendicecodigo
\apendicecodigotrue

\fancypagestyle{textual}{%
    \fancyhf{}%
    \fancyhead[L]{\small\textsf{\nouppercase{\leftmark}}}%
    \fancyhead[R]{\small\textsf{\textbf{\thepage}}}%
    \fancyfoot{}%
    \renewcommand{\headrulewidth}{0.4pt}%
    \renewcommand{\footrulewidth}{0pt}%
}

\fancypagestyle{plain}{%
    \fancyhf{}%
    \fancyhead[R]{\small\textsf{\textbf{\thepage}}}%
    \fancyfoot{}%
    \renewcommand{\headrulewidth}{0pt}%
    \renewcommand{\footrulewidth}{0pt}%
}

\usepackage{titlesec}
\usepackage{setspace}
\usepackage[plainpages=false,pdfpagelabels]{hyperref}

\hypersetup{
    colorlinks=true,
    linkcolor=blue!80!black,
    citecolor=green!50!black,
    urlcolor=blue!60!black
}

% =========================================================================
% LISTAGENS DE CÓDIGO-FONTE C++ (LISTINGS)
% =========================================================================

% -------------------------------------------------------------------------
% Paleta de Cores de Sintaxe
% -------------------------------------------------------------------------
\definecolor{codegreen}{rgb}{0,0.6,0}
\definecolor{codegray}{rgb}{0.5,0.5,0.5}
\definecolor{codepurple}{rgb}{0.58,0,0.82}
\definecolor{backcolour}{rgb}{0.9,0.9,0.9}

% -------------------------------------------------------------------------
% Configurações Globais do Listings
% -------------------------------------------------------------------------
\lstset{
    language=C++,
    basicstyle=\ttfamily\small,
    tabsize=4,
    keepspaces=true,
    columns=fullflexible,
    lineskip=-2pt,
    backgroundcolor=\color{backcolour},
    commentstyle=\color{codegreen},
    keywordstyle=\color{blue}\bfseries,
    stringstyle=\color{codepurple},
    deletestring={[b]'},
    numbers=left,
    numberstyle=\tiny\color{codegray},
    numbersep=6pt,
    breaklines=true,
    breakatwhitespace=false,
    showspaces=false,
    showstringspaces=false,
    showtabs=false,
    frame=single,
    rulecolor=\color{black},
    captionpos=b
}

% =========================================================================
% ESTRUTURAS SEMÂNTICAS, TEOREMAS E DEFINIÇÕES
% =========================================================================
\newtheorem{prop}{Proposição}

% -------------------------------------------------------------------------
% Definições Formais
% -------------------------------------------------------------------------
\newcounter{definicao}
\renewcommand{\thedefinicao}{\arabic{definicao}}
\NewDocumentEnvironment{definicao}{o o}{%
    \par\addvspace{0.8\baselineskip}\noindent
    \refstepcounter{definicao}%
    \IfValueT{#2}{\label{#2}}%
    \textbf{Definição~\thedefinicao\IfValueT{#1}{ (#1)}}.%
    \space
}{%
    \par\addvspace{0.5\baselineskip}
}

% -------------------------------------------------------------------------
% Teoremas e Demonstrações
% -------------------------------------------------------------------------
\newcounter{teorema}
\renewcommand{\theteorema}{\arabic{teorema}}
\NewDocumentEnvironment{teorema}{o o}{%
    \par\addvspace{0.8\baselineskip}\noindent
    \refstepcounter{teorema}%
    \IfValueT{#2}{\label{#2}}%
    \textbf{Teorema~\theteorema\IfValueT{#1}{ (#1)}}.%
    \space\itshape
}{%
    \par\addvspace{0.5\baselineskip}
}

\NewDocumentEnvironment{demonstracao}{}{%
    \par\noindent\textbf{Demonstração.}\space
}{%
    \hfill $\blacksquare$\par\addvspace{0.5\baselineskip}
}

% -------------------------------------------------------------------------
% Macros Conceituais e Auxiliares
% -------------------------------------------------------------------------
\newcommand{\intuicao}[1]{\paragraph{Intuição #1.}}
\newcommand{\implicacoes}[1]{\paragraph{Implicações Estratégicas #1.}}
\newcommand{\sa}{\text{s.a.}\quad}

% =========================================================================
% PALETA CROMÁTICA SEMÂNTICA (DESIGN SYSTEM)
% =========================================================================

% -------------------------------------------------------------------------
% Partições S e T
% -------------------------------------------------------------------------
\definecolor{setSborder}{RGB}{22, 101, 52}
\definecolor{setStext}{RGB}{20, 83, 45}
\definecolor{setSsource}{RGB}{187, 247, 208}
\definecolor{setSnode}{RGB}{220, 252, 231}
\colorlet{setSbg}{green!10}

\definecolor{setTborder}{RGB}{190, 18, 60}
\definecolor{setTtext}{RGB}{159, 18, 57}
\definecolor{setTsink}{RGB}{254, 205, 211}
\definecolor{setTnode}{RGB}{255, 228, 230}
\colorlet{setTbg}{red!10}

% -------------------------------------------------------------------------
% Fronteira de Corte e Elementos Neutros
% -------------------------------------------------------------------------
\definecolor{cutviolet}{RGB}{109, 40, 217}
\definecolor{cutred}{RGB}{225, 29, 72}
\definecolor{darkslate}{RGB}{30, 41, 59}

% -------------------------------------------------------------------------
% Cores da Capa e Cards
% -------------------------------------------------------------------------
\definecolor{topbg}{RGB}{15, 23, 42}
\definecolor{bottombg}{RGB}{248, 250, 252}
\definecolor{accentgold}{RGB}{226, 192, 116}
\definecolor{accentcyan}{RGB}{14, 165, 233}
\definecolor{cardbg}{RGB}{224, 232, 242}
\definecolor{cardborder}{RGB}{186, 202, 220}

% -------------------------------------------------------------------------
% Aliases de Compatibilidade
% -------------------------------------------------------------------------
\colorlet{scyan}{setSborder}
\colorlet{scyanbg}{setSnode}
\colorlet{ssourcebg}{setSsource}
\colorlet{tgold}{setTborder}
\colorlet{tgoldbg}{setTnode}
\colorlet{tsinkbg}{setTsink}

% =========================================================================
% BIBLIOTECA DE ESTILOS GRÁFICOS (TIKZ)
% =========================================================================
\newcommand{\tikzsubcaption}[2]{\node[font=\footnotesize\bfseries, above] at #1 {#2};}

\tikzset{
    >=Stealth,
    % ---------------------------------------------------------------------
    % Redes de Fluxo e Otimização
    % ---------------------------------------------------------------------
    flow node/.style={circle, draw, fill=white, minimum size=8.5mm, inner sep=1pt, font=\small},
    source node/.style={flow node, draw=setSborder, fill=setSsource, text=setStext, line width=1.1pt},
    sink node/.style={flow node, draw=setTborder, fill=setTsink, text=setTtext, line width=1.1pt},
    s node internal/.style={flow node, draw=setSborder!80, fill=setSnode, text=setStext, line width=0.9pt},
    t node internal/.style={flow node, draw=setTborder!80, fill=setTnode, text=setTtext, line width=0.9pt},
    trans node/.style={flow node, fill=yellow!15},
    flow edge/.style={->, thick},
    sat edge/.style={->, line width=1.5pt, cutviolet},
    res edge/.style={->, blue, dashed},
    rev edge/.style={->, red, dashed, line width=1.1pt},
    cut S/.style={dashed, draw=setSborder, fill=setSbg, rounded corners, thick},
    cut T/.style={dashed, draw=setTborder, fill=setTbg, rounded corners, thick},
    % ---------------------------------------------------------------------
    % Rótulos de Partição e Legendas
    % ---------------------------------------------------------------------
    cut label S/.style={font=\footnotesize\bfseries, text=setStext, anchor=north west},
    cut label T/.style={font=\footnotesize\bfseries, text=setTtext, anchor=north east},
    bip label L/.style={font=\large\bfseries, blue!80!black},
    bip label R/.style={font=\large\bfseries, orange!80!black},
    % ---------------------------------------------------------------------
    % Network Simplex e Árvores de Base
    % ---------------------------------------------------------------------
    tree edge/.style={->, line width=1.8pt, blue!80!black},
    nontree edge/.style={->, gray, dashed},
    pivot enter/.style={->, red, dashed, line width=1.5pt},
    pivot leave/.style={->, gray, dashed, line width=1.2pt},
    % ---------------------------------------------------------------------
    % Grafos em Camadas e Distâncias
    % ---------------------------------------------------------------------
    active node/.style={flow node, fill=yellow!30},
    gap node/.style={flow node, fill=orange!40},
    layer line/.style={gray!40, dashed},
    % ---------------------------------------------------------------------
    % Grafos Bipartidos e Emparelhamento
    % ---------------------------------------------------------------------
    bip L/.style={circle, draw, fill=blue!15, minimum size=7.5mm, inner sep=1pt, font=\small},
    bip R/.style={circle, draw, fill=orange!15, minimum size=7.5mm, inner sep=1pt, font=\small},
    cover node/.style={double, double distance=1.2pt},
    match edge/.style={-, line width=1.8pt, blue!80!black},
    unmatch edge/.style={-, gray},
    % ---------------------------------------------------------------------
    % Segmentação de Imagens e Caixas Conceituais
    % ---------------------------------------------------------------------
    pixel node/.style={circle, draw, fill=white, minimum size=8mm, font=\footnotesize},
    nlink/.style={-, thick},
    tlink s/.style={->, green!60!black, dashed},
    tlink t/.style={->, red!70!black, dashed},
    concept box/.style={draw, rounded corners=4pt, minimum height=1.4cm, align=center, font=\small},
    legend node/.style={circle, draw, minimum size=4mm, inner sep=0pt},
    % ---------------------------------------------------------------------
    % Rótulos de Arestas
    % ---------------------------------------------------------------------
    edge lbl/.style={font=\footnotesize, fill=white, inner sep=1.2pt, pos=0.5},
    edge lbl sloped/.style={edge lbl, sloped},
    edge lbl above/.style={edge lbl sloped, above=2.8pt},
    edge lbl below/.style={edge lbl sloped, below=2.8pt},
    edge lbl right/.style={edge lbl, right=3pt},
    edge lbl left/.style={edge lbl, left=3pt},
    edge lbl above norm/.style={edge lbl, above=3pt},
    edge lbl below norm/.style={edge lbl, below=3pt},
    edge lbl green/.style={font=\footnotesize, fill=green!10, inner sep=1.2pt, pos=0.5},
    edge lbl red/.style={font=\footnotesize, fill=red!10, inner sep=1.2pt, pos=0.5},
    stage box/.style={draw=blue!80!black, fill=blue!5, rounded corners=4pt, line width=1pt, align=center, font=\footnotesize, minimum height=1.4cm, text width=2.4cm},
    main/.style={flow node},
    % ---------------------------------------------------------------------
    % Pílulas da Capa e Badges
    % ---------------------------------------------------------------------
    cover badge/.style={
        draw,
        rounded corners=3pt,
        fill=topbg!95,
        font=\fontsize{7.5}{9.5}\selectfont\bfseries,
        minimum height=0.48cm,
        text height=5.2pt,
        text depth=0pt,
        inner xsep=6pt,
        inner ysep=0pt,
        align=center
    },
    card badge/.style={
        draw,
        rounded corners=3.5pt,
        fill=white,
        font=\fontsize{7.5}{9.5}\selectfont\bfseries,
        minimum height=0.52cm,
        text height=5.8pt,
        text depth=1.2pt,
        inner xsep=8pt,
        inner ysep=0pt,
        align=center
    }
}

% =========================================================================
% OPERADORES MATEMÁTICOS E MACROS AUXILIARES
% =========================================================================
\DeclareMathOperator{\ccap}{cap}
\DeclareMathOperator{\cost}{custo}
\DeclareMathOperator{\dist}{dist}

\newcommand{\cred}[2][]{\bar{w}_{#1}(#2)}
\newcommand{\cres}[2][f]{c_{#1}(#2)}
\newcommand{\Dres}[1][f]{D_{#1}}
\newcommand{\code}[1]{\lstinline{#1}}
\newcommand{\bannertext}[1]{\textbf{\textsf{#1}}}

\begin{document}
\setlength{\emergencystretch}{3em}
\onehalfspacing

% =========================================================================
% CAPA INSTITUCIONAL UFABC E TERMO DE HOMOLOGAÇÃO
% =========================================================================
\thispagestyle{empty}

\begin{center}
    {\large\textbf{FUNDAÇÃO UNIVERSIDADE FEDERAL DO ABC}}\\[0.15cm]
    {\large\textbf{CENTRO DE MATEMÁTICA, COMPUTAÇÃO E COGNIÇÃO}}\\[0.15cm]
    {\textbf{PRÓ-REITORIA DE PESQUISA --- PROPES}}\\[0.1cm]
    {\small\textbf{PROGRAMA INSTITUCIONAL DE INICIAÇÃO CIENTÍFICA (PIBIC / CNPq)}}\\[0.1cm]
    {\footnotesize Edital @@EDITAL@@}
\end{center}

\vspace{1.0cm}

\begin{center}
    {\LARGE\textbf{@@RESEARCH-TITLE@@:}}\\[0.3cm]
    {\Large\textbf{@@RESEARCH-SUBTITLE@@}}\par
\end{center}

\vspace{0.8cm}

\begin{flushright}
\begin{minipage}{0.6\textwidth}
\small
\setstretch{1.15}
\textbf{Relatório Final de Iniciação Científica} apresentado à Pró-Reitoria de Pesquisa da Fundação Universidade Federal do ABC (UFABC), como requisito de conclusão das atividades de pesquisa do Edital @@EDITAL@@.\par
\vspace{0.3cm}
\textbf{Pesquisador:} @@AUTHOR@@\\
\textbf{Orientação:} @@ADVISOR@@\\
\textbf{Área de Conhecimento:} Ciência da Computação / Matemática Discreta
\end{minipage}
\end{flushright}

\vspace{0.8cm}

\begin{center}
\small
\setstretch{1.15}
\begin{tabular}{|p{0.32\textwidth}|>{\raggedright\arraybackslash}p{0.60\textwidth}|}
\hline
\textbf{Título do Projeto:} & @@RESEARCH-TITLE@@: @@RESEARCH-SUBTITLE@@ \\ \hline
\textbf{Edital Vinculado:} & Edital @@EDITAL@@ \\ \hline
\textbf{Discente Pesquisador:} & @@AUTHOR@@ \\ \hline
\textbf{Docente Orientação:} & @@ADVISOR@@ \\ \hline
\textbf{Palavras-chave:} & Otimização Combinatória, Algoritmos, C++23 \\ \hline
\end{tabular}
\end{center}

\vfill

\begin{center}
\small
\begin{tabular}{cp{1.5cm}c}
\rule{6.5cm}{0.4pt} & & \rule{6.5cm}{0.4pt}\\
\textbf{@@AUTHOR@@} & & \textbf{@@ADVISOR@@}\\
Discente Pesquisador & & Orientação Acadêmica
\end{tabular}
\end{center}

\vspace{0.8cm}

\begin{center}
\footnotesize
Santo André --- SP\\
@@YEAR@@
\end{center}

\newpage

% =========================================================================
% SUMÁRIO (ANTES DO RESUMO)
% =========================================================================
\pagestyle{empty}
\renewcommand{\contentsname}{Sumário}
\setcounter{tocdepth}{1}
\tableofcontents

\clearpage

% =========================================================================
% RESUMO E INTRODUÇÃO
% =========================================================================
\pagestyle{textual}

% =========================================================================
% SEÇÃO 1: RESUMO
% =========================================================================
\section{Resumo}\label{sec:resumo}

@@RESEARCH-DESC@@ O estudo integra fundamentação matemática, implementação modular em C++23 e validação empírica reprodutível~\cite{cormen2022introduction, korte2018combinatorial}.

\vspace{0.35cm}

% =========================================================================
% SEÇÃO 2: INTRODUÇÃO
% =========================================================================
\section{Introdução}\label{sec:introducao}

Contextualize o problema de pesquisa @@RESEARCH-TITLE@@, suas propriedades matemáticas centrais e as motivações para o desenvolvimento algorítmico em C++23~\cite{cormen2022introduction, korte2018combinatorial}.

\clearpage

% =========================================================================
% METODOLOGIA
% =========================================================================
% =========================================================================
% SEÇÃO 3: METODOLOGIA
% =========================================================================
\section{Metodologia}\label{sec:metodologia}

\subsection{Materiais e Métodos}
Os algoritmos são implementados em C++23 sob interface base abstrata (\lstinline{AlgorithmBase}) e avaliados com medição isolada de tempo de execução via \lstinline{std::chrono::steady_clock}.

\subsection{Etapas da Pesquisa}
\begin{itemize}[leftmargin=1.5em, itemsep=0.1em, topsep=0.15em, parsep=0em]
    \item \textbf{Etapa 1:} Estudo teórico e formalização matemática do domínio.
    \item \textbf{Etapa 2:} Implementação modular em C++23 e suíte de testes automatizados.
    \item \textbf{Etapa 3:} Validação experimental e consolidação dos resultados.
\end{itemize}

\clearpage

% =========================================================================
% FUNDAMENTAÇÃO, RESULTADOS E APLICAÇÃO
% =========================================================================
% =========================================================================
% SEÇÃO 4: FUNDAMENTAÇÃO TEÓRICA E IMPLEMENTAÇÕES
% =========================================================================
\section{Fundamentação Teórica e Algoritmos}\label{sec:fundamentacao}

Apresente as definições formais, teoremas estruturais e provas de corretude e complexidade que fundamentam as implementações em C++23.

\begin{definicao}[Problema Base][def:ic_base]
Defina formalmente a instância do problema estudado.
\end{definicao}

\begin{figure}[!ht]
\centering
\begin{tikzpicture}[thick]
    \node[source node] (s) at (0, 0) {$s$};
    \node[flow node] (u) at (3, 1.2) {$u$};
    \node[flow node] (v) at (3, -1.2) {$v$};
    \node[sink node] (t) at (6, 0) {$t$};

    \draw[flow edge] (s) -- node[edge lbl above] {$w_1$} (u);
    \draw[flow edge] (s) -- node[edge lbl below] {$w_2$} (v);
    \draw[flow edge] (u) -- node[edge lbl right] {$w_3$} (v);
    \draw[flow edge] (u) -- node[edge lbl above] {$w_4$} (t);
    \draw[flow edge] (v) -- node[edge lbl below] {$w_5$} (t);
\end{tikzpicture}
\caption{Representação esquemática de referência.}
\label{fig:ic_exemplo}
\end{figure}

% =========================================================================
% SEÇÃO 5: RESULTADOS EXPERIMENTAIS
% =========================================================================
\section{Resultados Experimentais}\label{sec:resultados}

Síntese dos experimentos executados com a suíte da IC (\code{make bench-ic}), comparando tempo médio de execução e validação cruzada das soluções.

\IfFileExists{tables/table_results.tex}{
\begin{table}[!ht]
\centering
\small
\begin{tabular}{lrrr}
\toprule
\textbf{Instância} & \textbf{Tamanho} & \textbf{Algoritmo} & \textbf{Tempo Médio (ms)} \\
\midrule
\texttt{sample\_01} & 1.000 & \texttt{SampleAlgorithm} & 0.001 \\
\bottomrule
\end{tabular}
\caption{Resultados experimentais iniciais da suíte da IC (gerados via \texttt{make bench-ic}).}
\label{tab:results_ic}
\end{table}
}{}

% =========================================================================
% SEÇÃO 6: APLICAÇÃO PRÁTICA
% =========================================================================
\section{Aplicação Prática}\label{sec:aplicacao}

Descreva a modelagem de uma aplicação real que utiliza os algoritmos desenvolvidos em \code{applications/}.

% =========================================================================
% SEÇÃO 7: CONCLUSÕES
% =========================================================================
\section{Conclusões}\label{sec:conclusoes}

Consolide os resultados teóricos, computacionais e experimentais alcançados na pesquisa \textbf{@@RESEARCH-TITLE@@}, além dos horizontes futuros na versão monográfica em livro.

\newpage

\section{Bibliografia}
\bibliographystyle{plain}
\bibliography{common/references}

\newpage

\appendix
\renewcommand{\theHsection}{apendice.\arabic{section}}
\providecommand{\apendice}[1]{\section{#1}}

% =========================================================================
% APÊNDICE: RESULTADOS EXPERIMENTAIS EXAUSTIVOS
% =========================================================================
\apendice{Resultados Experimentais Exaustivos}
\label{ap:benchmarks_completos}

Complementa a Seção~\ref{sec:resultados} com as medições completas da suíte experimental da Iniciação Científica.

\newpage

\begin{center}{\Large\textbf{Apêndice Computacional: Implementações em C++23}}\end{center}

% =========================================================================
% ARQUITETURA DE SOFTWARE EM C++23 E ESTRUTURAS DE DADOS
% =========================================================================
\apendice{Arquitetura Orientada a Objetos em C++23 e Estruturas de Dados}
\label{ap:arquitetura}

A biblioteca de algoritmos de \textbf{@@RESEARCH-TITLE@@} segue o padrão ISO C++23 sob orientação a objetos com os padrões \textit{Factory Method} e \textit{Prototype}~\cite{gamma1994design}.

\paragraph{Tipagem e Margem de Segurança Numérica.}
Adotam-se os apelidos padronizados \lstinline{Long} (\lstinline{long long}) e \lstinline{Size} (\lstinline{std::size_t}), com constantes de saturação definidas com deslocamento defensivo:
\begin{lstlisting}[language=C++, numbers=none, basicstyle=\ttfamily\footnotesize, frame=none, backgroundcolor=\color{white}]
constexpr Long INF = std::numeric_limits<Long>::max() >> 8;
constexpr Size MAX = std::numeric_limits<Size>::max() >> 8;
\end{lstlisting}

\paragraph{Hierarquia Polimórfica.}
Os algoritmos concretos herdam da classe base abstrata \lstinline{AlgorithmBase} (Apêndice~\ref{ap:base}), expondo os métodos virtuais \lstinline{make(n)}, \lstinline{clone()} e \lstinline{solve()}.

\newpage

% =========================================================================
% IMPLEMENTAÇÃO DA INTERFACE BASE ``ALGORITHMBASE''
% =========================================================================
\apendice{Implementação da Interface Base ``AlgorithmBase''}
\label{ap:base}

\begin{lstlisting}[caption={Implementação em C++23 da interface base abstrata \lstinline{AlgorithmBase}.}, label={lst:base}]
#ifndef ALGORITHM_BASE_HPP
#define ALGORITHM_BASE_HPP

#include <cstddef>
#include <limits>
#include <memory>
#include <vector>

using Long = long long;
using Size = std::size_t;

constexpr Long INF = std::numeric_limits<Long>::max() >> 8;
constexpr Size MAX = std::numeric_limits<Size>::max() >> 8;

class AlgorithmBase
{
public:
	explicit AlgorithmBase(const Size n) : size(n), data(n, 0) {}
	virtual ~AlgorithmBase() = default;

	virtual std::unique_ptr<AlgorithmBase> make(const Size n) const = 0;
	virtual std::unique_ptr<AlgorithmBase> clone() const = 0;

	virtual void add_item(const Size index, const Long value)
	{
		if (index < size)
			data[index] += value;
	}

	virtual Long solve() = 0;

	[[nodiscard]] Size get_size() const
	{
		return size;
	}

	[[nodiscard]] const std::vector<Long> &get_data() const
	{
		return data;
	}

protected:
	Size size;
	std::vector<Long> data;
};

#endif
\end{lstlisting}

\newpage

% =========================================================================
% IMPLEMENTAÇÃO DA CLASSE ``SAMPLEALGORITHM''
% =========================================================================
\apendice{Implementação da Classe ``SampleAlgorithm''}
\label{ap:sample_algorithm}

\begin{lstlisting}[caption={Implementação de referência \lstinline{SampleAlgorithm} em C++23.}, label={lst:sample_algorithm}]
#ifndef SAMPLE_ALGORITHM_HPP
#define SAMPLE_ALGORITHM_HPP

#include "../core/AlgorithmBase.hpp"

class SampleAlgorithm final : public AlgorithmBase
{
public:
	explicit SampleAlgorithm(const Size n) : AlgorithmBase(n) {}

	static std::unique_ptr<AlgorithmBase> create(const Size n)
	{
		return std::make_unique<SampleAlgorithm>(n);
	}

	std::unique_ptr<AlgorithmBase> make(const Size n) const override
	{
		return std::make_unique<SampleAlgorithm>(n);
	}

	std::unique_ptr<AlgorithmBase> clone() const override
	{
		return std::make_unique<SampleAlgorithm>(*this);
	}

	Long solve() override
	{
		Long total = 0;
		for (const Long value : data)
			total += value;
		return total;
	}
};

#endif
\end{lstlisting}

\newpage

% =========================================================================
% ACESSO AO CÓDIGO-FONTE E REPOSITÓRIO GITHUB
% =========================================================================
\apendice{Acesso ao Código-Fonte e Repositório GitHub}
\label{ap:github}

Para garantir a reprodutibilidade dos experimentos e disponibilizar o código-fonte desenvolvido, todo o material da pesquisa \textbf{@@RESEARCH-TITLE@@} encontra-se organizado em repositório público no GitHub:

\vspace{0.4cm}
\begin{center}
    \fbox{\url{https://github.com/@@GITHUB-USER@@/@@RESEARCH-SLUG@@}}
\end{center}
\vspace{0.4cm}

\end{document}
